\chapter{Stochastic orders, quadratic variation, and Taylor}
\label{chap:pre_orders}

This chapter collects the supporting API that the asymptotic proofs use but that
Mathlib does not yet package: a calculus for the stochastic Landau symbols
$o_p, O_p$; the predictable quadratic variation of a discrete-time martingale;
and a multivariate first-order Taylor expansion with a quadratic remainder. None
of these is a deep theorem, but each is used pervasively.

\section{The \texorpdfstring{$o_p$}{op} / \texorpdfstring{$O_p$}{Op} calculus}
\label{sec:pre_op}

Recall Definition~\ref{def:op_Op}. We reduce everything to statements about
$o_p(1)$ (convergence to $0$ in probability) and $O_p(1)$ (boundedness in
probability, i.e.\ tightness).

\textbf{Formalization note.} Mathlib provides the primitives: convergence in
probability is \texttt{MeasureTheory.TendstoInMeasure} to a constant
(\texttt{MeasureTheory/Function/ConvergenceInMeasure.lean}), and boundedness in
probability is tightness of the laws. The lemmas below are an API layer on top of
these, plus the bridge to Slutsky's theorem
(\texttt{ConvergenceInDistribution.lean}). This is engineering rather than new
mathematics, but it must exist before any of the rate computations can be
written.

\begin{lemma}[Reduction to order one]\label{lem:op_factor}
\uses{def:op_Op}
$X_n = o_p(a_n)$ if and only if $X_n/a_n = o_p(1)$, and $X_n = O_p(a_n)$ if and
only if $X_n/a_n = O_p(1)$.
\end{lemma}

\begin{proof}
Unfold Definition~\ref{def:op_Op}: both sides are the same statement about the
sequence $|X_n|/a_n$.
\end{proof}

\begin{lemma}[Convergence in probability gives $o_p$]\label{lem:op_of_tendsto}
\lean{AlphaRAR.isLittleOpOne_of_tendsto_ae}
\uses{def:op_Op}
If $X_n \to 0$ in probability (in particular, if $X_n \to 0$ a.s.), then
$X_n = o_p(1)$. If $X_n \to c$ in distribution for some random variable, then
$X_n = O_p(1)$.
\end{lemma}

\begin{proof}
The first statement is the definition of convergence in probability. For the
second, convergence in distribution implies tightness of the laws
($\{X_n\}$ is uniformly tight), which is exactly $O_p(1)$.
\end{proof}

\begin{lemma}[Arithmetic of stochastic orders]\label{lem:op_arith}
\lean{AlphaRAR.IsLittleOpOne.add, AlphaRAR.IsBigOpOne.add, AlphaRAR.IsLittleOpOne.add_bigOpOne, AlphaRAR.IsBigOpOne.mul, AlphaRAR.IsBigOpOne.mul_littleOp, AlphaRAR.IsBigOpOne.bdd_mul}\leanok
\uses{def:op_Op}
For a fixed rate $a_n$:
\begin{enumerate}
  \item[(i)] $o_p(a_n) + o_p(a_n) = o_p(a_n)$ and $O_p(a_n) + O_p(a_n) = O_p(a_n)$;
  \item[(ii)] $o_p(a_n) + O_p(a_n) = O_p(a_n)$;
  \item[(iii)] if $X_n = O_p(1)$ and $Y_n = o_p(1)$ then $X_n Y_n = o_p(1)$;
  \item[(iv)] if $X_n = O_p(a_n)$ and $c_n \to c \in \R$ then $c_n X_n = O_p(a_n)$.
\end{enumerate}
\end{lemma}

\begin{proof}\leanok
\uses{lem:op_factor}
By Lemma~\ref{lem:op_factor} reduce to $a_n = 1$. Then (i)--(ii) are union bounds
on the events $\{|X_n| > \varepsilon\}$, $\{|Y_n| > \varepsilon\}$; (iii) follows
by conditioning on the tightness of $X_n$: for $\eta$ pick $M$ with
$\sup_n \Prob(|X_n|>M)\le\eta$, then $\Prob(|X_nY_n|>\varepsilon)\le\eta
+\Prob(|Y_n|>\varepsilon/M)$; (iv) is immediate from boundedness of $c_n$.
\end{proof}

\begin{lemma}[Slutsky bridge]\label{lem:op_slutsky}
\uses{def:op_Op}
If $X_n \eqdist X$ (convergence in distribution) and $R_n = o_p(1)$, then
$X_n + R_n \eqdist X$. More generally, if $(X_n) \eqdist X$ and $Y_n \to c$ in
probability for a constant $c$, then $(X_n, Y_n) \eqdist (X, c)$.
\end{lemma}

\begin{proof}
\uses{lem:op_of_tendsto}
This is Slutsky's theorem, available in Mathlib as
\texttt{TendstoInDistribution.prodMk\_of\_tendstoInMeasure\_const}, combined with
the continuous mapping theorem \texttt{TendstoInDistribution.continuous\_comp}
applied to addition.
\end{proof}

\section{Predictable quadratic variation}
\label{sec:pre_qv}

\begin{definition}[Predictable quadratic variation]\label{def:pred_qv}
\lean{AlphaRAR.predQuadVar}\leanok
Let $(M_n)_{n\ge0}$ be a square-integrable $(\filt_n)$-martingale with $M_0 = 0$.
Its \emph{predictable quadratic variation} is
$\langle M\rangle := \mathrm{predictablePart}(M^2)$, the predictable part of the
submartingale $M^2$ in the Doob decomposition. Equivalently,
$\langle M\rangle_0 = 0$ and
$\langle M\rangle_n = \sum_{i=1}^n \E[(\Delta M_i)^2 \mid \filt_{i-1}]$.
\end{definition}

\textbf{Formalization note.} Mathlib already has the discrete Doob decomposition
(\texttt{predictablePart} / \texttt{martingalePart} in
\texttt{Probability/Martingale/Centering.lean}); $\langle M\rangle$ is a thin
wrapper. The lemmas below are the basic API and are straightforward from that
decomposition.

\begin{lemma}[Increment of the quadratic variation]\label{lem:qv_incr}
\lean{AlphaRAR.predQuadVar_succ_sub_eq, AlphaRAR.predQuadVar_ae_eq_sum, AlphaRAR.predQuadVar_mono,
AlphaRAR.isStronglyPredictable_predQuadVar}\leanok
\uses{def:pred_qv}
For a square-integrable martingale $M$ with $M_0=0$,
$\langle M\rangle_n - \langle M\rangle_{n-1} = \E[(\Delta M_n)^2\mid\filt_{n-1}]$,
and $\langle M\rangle$ is nondecreasing and predictable.
\end{lemma}

\begin{proof}\leanok
The Doob decomposition of the submartingale $M^2$ has predictable part with
increments $\E[M_n^2 - M_{n-1}^2\mid\filt_{n-1}]
= \E[(\Delta M_n)^2\mid\filt_{n-1}]$, using that $M$ is a martingale so the cross
term $\E[2M_{n-1}\Delta M_n\mid\filt_{n-1}]=0$. Nonnegativity of the increments
gives monotonicity.
\end{proof}

\begin{lemma}[Linear bound on $\langle M\rangle$ for bounded increments]\label{lem:qv_linear_bound}
\lean{AlphaRAR.predQuadVar_le_of_bound}\leanok
\uses{def:pred_qv, lem:qv_incr}
If $|\Delta M_i|\le c$ a.s., then $\langle M\rangle_n \le c^2 n$ a.s.
\end{lemma}

\begin{proof}\leanok
Each increment satisfies $\langle M\rangle_{k+1}-\langle M\rangle_k
= \E[(\Delta M_k)^2\mid\filt_k] \le \E[c^2\mid\filt_k] = c^2$; telescoping from
$\langle M\rangle_0 = 0$ gives $\langle M\rangle_n \le c^2 n$.
\end{proof}

\begin{lemma}[$M^2 - \langle M\rangle$ is a martingale]\label{lem:qv_mart}
\lean{AlphaRAR.martingale_sq_sub_predQuadVar}\leanok
\uses{def:pred_qv}
$M^2 - \langle M\rangle$ is an $(\filt_n)$-martingale.
\end{lemma}

\begin{proof}\leanok
This is the martingale part in the Doob decomposition of $M^2$
(\texttt{martingale\_martingalePart}).
\end{proof}

\begin{lemma}[$\langle M\rangle$ is the unique compensator]\label{lem:qv_unique}
\lean{AlphaRAR.IsPredQuadVar, AlphaRAR.isPredQuadVar_predQuadVar,
AlphaRAR.IsPredQuadVar.indistinguishable_predQuadVar}\leanok
\uses{def:pred_qv, lem:qv_mart}
Call a process $A$ a \emph{compensator} of $M^2$ if $M^2 - A$ is an $(\filt_n)$-martingale, $A$ is
predictable, $A_0 = 0$, and each $A_n$ is integrable. Then $\langle M\rangle$ is a compensator of
$M^2$, and any compensator $A$ of $M^2$ is indistinguishable from it:
$\P[\forall n,\ A_n = \langle M\rangle_n] = 1$. So Definition~\ref{def:pred_qv}, which is a
formula, defines the object the name describes, and does so up to the a.s.\ equality of paths that
is the best any statement about conditional expectations can offer.
\end{lemma}

\begin{proof}\leanok
\uses{lem:qv_mart, lem:qv_incr}
Existence is Lemma~\ref{lem:qv_mart} together with the predictability of the Doob predictable part
(Lemma~\ref{lem:qv_incr}). For uniqueness, $M^2 = (M^2 - A) + A$ exhibits $M^2$ as a martingale
plus a predictable process null at $0$, and the Doob decomposition is unique up to null sets
(\texttt{predictablePart\_add\_ae\_eq}), so $A_n = \langle M\rangle_n$ a.s.\ for each fixed $n$;
the index set is countable, so the countably many exceptional sets union to a null set.
\end{proof}

\begin{lemma}[Expected quadratic variation equals second moment]\label{lem:qv_second_moment}
\lean{AlphaRAR.integral_sq_eq_integral_predQuadVar,
AlphaRAR.integral_sq_eq_sum_integral_increment_sq}\leanok
\uses{def:pred_qv, lem:qv_mart}
For a square-integrable martingale $M$ with $M_0 = 0$ and every $n$,
$\E[M_n^2] = \E[\langle M\rangle_n]$.
\end{lemma}

\begin{proof}\leanok
$M^2 - \langle M\rangle$ is a martingale (Lemma~\ref{lem:qv_mart}) starting at
$M_0^2 - \langle M\rangle_0 = 0$, so its expectation is constant equal to $0$;
that is $\E[M_n^2] - \E[\langle M\rangle_n] = 0$.
\end{proof}

\begin{lemma}[Martingale $L^2$ growth]\label{lem:mart_sq_growth}
\lean{AlphaRAR.integral_sq_le_of_increment_bound}\leanok
\uses{def:pred_qv, lem:qv_second_moment, lem:qv_incr}
Let $M$ be a square-integrable martingale with $M_0 = 0$. If every increment has
second moment $\E[(\Delta M_{n+1})^2] \le \sigma^2$, then $\E[M_n^2] \le \sigma^2 n$.
\end{lemma}

\begin{proof}\leanok
By Lemma~\ref{lem:qv_second_moment}, $\E[M_n^2] = \E[\langle M\rangle_n]$.
Telescoping and Lemma~\ref{lem:qv_incr},
$\E[\langle M\rangle_n] = \sum_{i=0}^{n-1}\E[(\Delta M_{i+1})^2] \le n\sigma^2$.
\end{proof}

\begin{lemma}[Orthogonal martingales]\label{lem:qv_orthogonal}
\lean{AlphaRAR.martingale_mul, AlphaRAR.predQuadVar_add_of_martingale_mul}\leanok
\uses{def:pred_qv}
Let $M, N$ be square-integrable martingales such that
$\E[\Delta M_i\, \Delta N_i \mid \filt_{i-1}] = 0$ for all $i$ (e.g.\ if
$\Delta M_i\, \Delta N_i = 0$ a.s.). Then $MN$ is a martingale, so the cross
variation vanishes, and $\langle M + N\rangle = \langle M\rangle + \langle N\rangle$.
\end{lemma}

\begin{proof}\leanok
$\E[\Delta(MN)_i\mid\filt_{i-1}] = M_{i-1}\E[\Delta N_i\mid\cdot]
+ N_{i-1}\E[\Delta M_i\mid\cdot] + \E[\Delta M_i\Delta N_i\mid\cdot] = 0$, so $MN$
is a martingale (\texttt{martingale\_mul}). For additivity, expand
$(M+N)^2 = M^2 + N^2 + 2MN$; the predictable part is linear and the predictable
part of the martingale $MN$ vanishes, so
$\langle M+N\rangle = \langle M\rangle + \langle N\rangle$.
\end{proof}

\section{A multivariate Taylor remainder}
\label{sec:pre_taylor}

\begin{lemma}[First-order expansion with quadratic remainder]\label{lem:taylor_first_order}
\lean{AlphaRAR.norm_sub_fderiv_le_mul_sq}\leanok
Let $f : \R^K \to \R^d$ be $C^2$ on a convex open set $\hilb$, and suppose
$\|\mathrm{D}^2 f(z)\| \le C$ for all $z$ in a convex neighborhood of $y \in \hilb$.
Then for $x$ in that neighborhood,
\[
  \big\| f(x) - f(y) - \mathrm{D}f(y)(x - y) \big\| \le \tfrac{C}{2}\,\|x - y\|^2.
\]
\end{lemma}

\begin{proof}\leanok
Apply the mean-value / Taylor formula to $t \mapsto f(y + t(x-y))$ on $[0,1]$:
$f(x) - f(y) - \mathrm{D}f(y)(x-y)
= \int_0^1 (1-t)\, \mathrm{D}^2 f(y + t(x-y))[x-y, x-y]\, dt$, and bound the
integrand by $C\|x-y\|^2$.

\emph{Formalization.} \texttt{AlphaRAR.norm\_sub\_fderiv\_le\_mul\_sq} states the equivalent bound
$\|f(x)-f(y)-\mathrm{D}f(y)(x-y)\| \le K\|x-y\|^2$ under the (weaker, and here more convenient)
hypothesis that the derivative is $K$-Lipschitz relative to $y$ on a convex set,
$\|\mathrm{D}f(z)-\mathrm{D}f(y)\| \le K\|z-y\|$ --- which the second-derivative bound
$\|\mathrm{D}^2 f\|\le C$ supplies with $K = C$ by the mean value inequality. The proof applies the
mean value inequality \texttt{Convex.norm\_image\_sub\_le\_of\_norm\_hasFDerivWithin\_le'} on the
segment $[y,x]$, where $\|\mathrm{D}f(z)-\mathrm{D}f(y)\|\le K\|z-y\|\le K\|x-y\|$ gives the constant.
\end{proof}

\textbf{Formalization note.} Mathlib has the pieces —
\texttt{fderiv}, \texttt{ContDiff}, iterated derivatives, and one-dimensional
Taylor with remainder (\texttt{Analysis/Calculus/Taylor.lean}). What is missing
is this packaged multivariate first-order statement with an explicit quadratic
bound on a neighborhood; it can be assembled from the existing bounds on the
second derivative and the integral form of the remainder.

\begin{corollary}[Expansion of the target function]\label{lem:taylor_rho}
\uses{def:condB, lem:taylor_first_order}
Under Condition~\textbf{B}, there is a constant $C$ and a neighborhood of
$\Params$ on which
$\big\|\Target(x) - \Target(\Params) - G(x - \Params)\big\| \le C\|x - \Params\|^2$,
where $G = \nabla\Target|_\Params$. In particular, if $\hat X \to \Params$ with
$\hat X - \Params = O_p(r_n)$, then
$\Target(\hat X) - \Target(\Params) = G(\hat X - \Params) + O_p(r_n^2)$.
\end{corollary}

\begin{proof}
\uses{lem:op_arith}
Condition~\textbf{B} makes $\Target$ twice differentiable near $\Params$, so
Lemma~\ref{lem:taylor_first_order} applies with $y = \Params$; the second
statement follows from the order arithmetic of Lemma~\ref{lem:op_arith}.
\end{proof}
